\begin{frame}{수용장: 뒤쪽 층이 ``보게 되는'' 원본 영역}
  풀링의 공간 축소를 다른 관점에서 보면 \kb{수용장(receptive field)}
  의 \textbf{확대}. 합성곱 한 층은 입력의 $k\times k$ 영역만
  직접 보지만, 출력을 \textbf{또} 합성곱·풀링하면 두 번째 층의
  뉴런 하나가 ``보게 되는'' 원본 영역은 더 넓어진다.
  \vskip0.4em
  각 층의 필터 크기 $k_l$, 스트라이드 $s_l$ 일 때, $l$번째 층
  뉴런의 수용장 크기:
  \[ R_l = R_{l-1} + (k_l - 1)\cdot \prod_{j < l} s_j,
     \qquad R_0 = 1 \]
  \begin{itemize}
    \item 한 층의 뉴런은 직전 층의 뉴런 $k_l$개를 ``보는데'',
      각 뉴런이 이미 원본의 $R_{l-1}$ 픽셀을 대표
    \item 인접 $k_l$개의 수용장을 합치면 겹치는 $(k_l-1)$ 픽셀만
      \textbf{새로} 보이고, 이 확장량이 이전 스트라이드 곱만큼
      \textbf{스케일업}
  \end{itemize}
\end{frame}
